% ============================================================
% Section 145: 掺磷硅片的表面态与表面势
% ============================================================
\section{半导体物理：掺磷硅片的表面态与表面势}
\label{sec:surface-states-silicon}

\begin{modelbox}[题目：表面态对硅表面势及能带的影响]
一块抛光硅片，体内掺磷 $N_D=1\times10^{15}\,\text{cm}^{-3}$，设有均匀分布的表面态，密度为 $1\times10^{13}\,\text{cm}^{-2}\cdot\text{eV}^{-1}$，且可用位于 $E_v+0.3\,\text{eV}$ 的中性能级表示。求室温下：

(1) 体内载流子浓度；

(2) 体内费米能级位置；

(3) 自由硅表面势；

(4) 画出能带图，注明能量关系，并计算说明原因。

（已知 $E_g=1.12\,\text{eV}$，$N_c=2.8\times10^{19}\,\text{cm}^{-3}$，$n_i=1.5\times10^{10}\,\text{cm}^{-3}$，$\varepsilon_r=11.8$）
\end{modelbox}

\begin{tipbox}[考点快速分析]
\begin{itemize}
    \item \textbf{体内参数}：施主全电离，$n_0\approx N_D$，$p_0=n_i^2/n_0$
    \item \textbf{费米能级}：$E_F=E_c-k_BT\ln(N_c/n_0)$
    \item \textbf{表面势}：由体内费米能级与表面中性能级之差决定
    \item \textbf{耗尽层}：电子流向表面，体内留下电离施主正电荷
\end{itemize}
\end{tipbox}

\begin{derivationbox}[标准解题过程]
\textbf{(1) 体内载流子浓度}

磷为施主杂质，室温下全电离：
\[
n_0\approx N_D=1\times10^{15}\,\text{cm}^{-3}.
\]
空穴浓度：
\[
p_0=\frac{n_i^2}{n_0}=\frac{2.25\times10^{20}}{10^{15}}=2.25\times10^5\,\text{cm}^{-3}.
\]

\textbf{(2) 体内费米能级位置}

由 $n_0=N_ce^{-(E_c-E_F)/k_BT}$：
\[
E_F=E_c-k_BT\ln\frac{N_c}{n_0}=E_c-0.026\ln(2.8\times10^4)\approx E_c-0.27\,\text{eV}.
\]
以价带顶 $E_v=0$ 为参考，$E_c=E_g=1.12\,\text{eV}$，故
\begin{equation}
\boxed{E_{F,\text{bulk}}=E_c-0.27=0.85\,\text{eV}.}
\end{equation}

\textbf{(3) 自由硅表面势}

表面中性能级 $E_t=E_v+0.3\,\text{eV}=0.3\,\text{eV}$。体内费米能级 $E_{F,\text{bulk}}=0.85\,\text{eV}>E_t$，电子从体内流向表面，表面带负电，体内形成正空间电荷区（耗尽层）。

平衡时表面与体内费米能级拉平，能带向上弯曲。表面势：
\begin{equation}
\boxed{eV_s=E_{F,\text{bulk}}-E_t=0.85-0.30=0.55\,\text{eV}\implies V_s=0.55\,\text{V}.}
\end{equation}

\textbf{(4) 能带图与空间电荷层参数}

耗尽层宽度（突变耗尽近似）：
\[
x_d=\sqrt{\frac{2\varepsilon_0\varepsilon_rV_s}{eN_D}}\approx\sqrt{\frac{2\times8.85\times10^{-14}\times11.8\times0.55}{1.6\times10^{-19}\times10^{15}}}\approx8.47\times10^{-5}\,\text{cm}=0.847\,\mu\text{m}.
\]

单位面积耗尽层电荷：
\[
Q_d=eN_Dx_d\approx1.6\times10^{-19}\times10^{15}\times8.47\times10^{-5}\approx1.36\times10^{-8}\,\text{C/cm}^2.
\]

表面态电荷变化：表面态密度 $D_{it}=10^{13}\,\text{cm}^{-2}\cdot\text{eV}^{-1}$。为平衡耗尽层正电荷，表面费米能级需移动 $\Delta E=(Q_d/e)/D_{it}\approx8.47\times10^{-3}\,\text{eV}$，此移动极小，表面中性能级几乎仍位于 $E_v+0.3\,\text{eV}$。
\end{derivationbox}

\begin{formulabox}[答案汇总]
\begin{center}
\begin{tabular}{ll}
\toprule
\textbf{项目} & \textbf{数值} \\
\midrule
(1) 体内电子浓度 $n_0$ & $1\times10^{15}\,\text{cm}^{-3}$ \\
(2) 体内费米能级 $E_F$ & $E_c-0.27\,\text{eV}$（或 $E_v+0.85\,\text{eV}$） \\
(3) 表面势 $V_s$ & $0.55\,\text{V}$ \\
(4) 耗尽层宽度 $x_d$ & $8.47\times10^{-5}\,\text{cm}=0.847\,\mu\text{m}$ \\
\bottomrule
\end{tabular}
\end{center}
\end{formulabox}

\begin{insightbox}[物理图像]
\begin{itemize}
    \item 表面态的存在导致能带弯曲，形成表面空间电荷层。当体内费米能级高于表面中性能级时，电子向表面转移，n型半导体表面形成耗尽层。
    \item 表面势 $V_s$ 由体内费米能级与表面中性能级的能量差决定。
    \item 表面态密度较高时，表面费米能级被``钉扎''在中性能级附近，外加偏压主要改变耗尽层宽度，而非表面势。
\end{itemize}
\end{insightbox}
